\documentclass[11pt]{article}
\usepackage[margin=1in]{geometry}
\usepackage{amsmath, amssymb, amsthm}
\usepackage{hyperref}
\usepackage{graphicx}
\usepackage{booktabs}

\title{A Fair Value Engine for Kalshi Sports Markets:\\
\large Binary Option Pricing via Brownian Motion}
\author{}
\date{\today}

\newtheorem{proposition}{Proposition}

\begin{document}
\maketitle

\begin{abstract}
Kalshi event contracts settle at \$1 if an event occurs and \$0 otherwise,
which makes them mathematically identical to cash-or-nothing binary options.
This document derives a fair value model for Kalshi ``team X wins'' game
markets in two layers: a pre-game Elo-based prior, and an in-game model
where the score margin is treated as a Brownian motion and the win
probability follows directly from the distribution of its terminal value.
We then use this fair value to define market edge and size positions via
the Kelly criterion. Every result here follows from undergraduate
probability: the CDF of the normal distribution, elementary properties of
Brownian motion, and Kelly's log-utility argument.
\end{abstract}

\section{Kalshi contracts as binary options}

A Kalshi ``YES'' contract on an event $E$ pays
\[
\text{Payoff} =
\begin{cases}
\$1 & \text{if } E \text{ occurs} \\
\$0 & \text{otherwise}
\end{cases}
\]
Under the assumption of no arbitrage and risk-neutral pricing, the fair
price of this contract today is
\[
V_0 = e^{-r\tau}\,\mathbb{E}^{\mathbb{Q}}\!\left[\mathbf{1}_E\right]
    = e^{-r\tau}\, \mathbb{Q}(E),
\]
where $\mathbb{Q}$ is the risk-neutral measure and $\tau$ is the (typically
short, hours-to-days) time to settlement. Because $\tau$ is small for a
single game, $e^{-r\tau}\approx 1$, so
\[
V_0 \approx \mathbb{Q}(E).
\]
This is the key fact that makes Kalshi different from a normal asset: the
\emph{price itself} is (approximately) a probability. A market quote of
63\textcent{} is a direct claim: ``the market believes there is a 63\%
chance of $E$.'' Our entire project is about producing an independent
estimate of $\mathbb{Q}(E)$ (call it $p_{\text{model}}$) and comparing it to
the market's implied estimate $p_{\text{market}} = V_0$.

\section{Layer 1: pre-game probability via Elo ratings}

Before a game starts, we need a prior belief about each team's strength.
The Elo system (originally designed for chess) gives each team $i$ a
rating $R_i \in \mathbb{R}$, and converts a rating difference into a win
probability via the logistic function:
\[
p_{\text{home}} = \frac{1}{1 + 10^{-(R_{\text{home}} + H - R_{\text{away}})/400}}
\]
where $H$ is a fixed home-field-advantage bonus (in rating points). After
each game, ratings update toward the observed outcome:
\[
R_{\text{home}}' = R_{\text{home}} + K\big(S_{\text{home}} - p_{\text{home}}\big),
\qquad S_{\text{home}} \in \{0, 1\}
\]
This is a simple stochastic approximation / gradient update: it nudges the
rating in the direction of the surprise $(S - p)$, scaled by learning rate
$K$. Elo can be derived formally as an online logistic regression with a
single feature (rating difference) updated by gradient descent on the log
loss, which is worth knowing but not required to use it. In this project
Elo supplies the \textbf{prior} at kickoff ($\tau = 1$ below); the in-game
model in Section 3 is what does the interesting stochastic work.

\section{Layer 2: in-game probability via Brownian motion (Stern 1994)}

\subsection{Model setup}

Let $M(t)$ denote the point margin (home score minus away score) at time
$t$ during the game, and let $\tau(t) \in (0, 1]$ denote the \emph{fraction
of the game remaining} at time $t$ (so $\tau = 1$ at kickoff and $\tau \to
0$ at the final whistle).

\textbf{Assumption.} The margin process behaves like a driftless Brownian
motion in ``fraction-of-game-remaining'' time:
\[
M(t) - M(t_0) \;\sim\; \mathcal{N}\!\big(0,\ \sigma^2 (\tau_0 - \tau)\big)
\]
for $t > t_0$, where $\sigma^2$ is a fixed volatility parameter (units:
points$^2$ per full game).

\textbf{Why this is a reasonable approximation.} Scoring in football and
basketball happens through a sequence of discrete, roughly independent
possessions, each contributing a small random point swing. By the
functional Central Limit Theorem (Donsker's theorem), a sum of many small,
roughly independent, roughly identically distributed increments converges
to a Brownian motion as the number of increments grows. A game has enough
possessions (order of 10--20 in the NFL, far more in the NBA) that this
approximation is reasonable, though it is not exact -- see the limitations
in Section 6.

\subsection{Deriving the win probability formula}

Fix the current time $t_0$ with current margin $m = M(t_0)$ and fraction of
game remaining $\tau = \tau(t_0)$. The home team wins if and only if the
\emph{final} margin $M_{\text{final}} = M(1) $ (using $\tau=0$ as the
endpoint) is positive.

By the model assumption, conditional on the current state,
\[
M_{\text{final}} \;\sim\; \mathcal{N}\!\big(m,\ \sigma^2 \tau\big).
\]

\begin{proposition}
Under the above model,
\[
\mathbb{P}(\text{home wins} \mid m, \tau) = \Phi\!\left(\frac{m}{\sigma\sqrt{\tau}}\right)
\]
where $\Phi$ is the standard normal CDF.
\end{proposition}

\begin{proof}
Standardize $M_{\text{final}}$: let $Z = \dfrac{M_{\text{final}} - m}{\sigma\sqrt{\tau}}$,
so $Z \sim \mathcal{N}(0,1)$. Then
\[
\mathbb{P}(M_{\text{final}} > 0)
= \mathbb{P}\!\left(Z > \frac{-m}{\sigma\sqrt{\tau}}\right)
= 1 - \Phi\!\left(\frac{-m}{\sigma\sqrt{\tau}}\right)
= \Phi\!\left(\frac{m}{\sigma\sqrt{\tau}}\right),
\]
using the symmetry of the standard normal distribution, $1-\Phi(-x) = \Phi(x)$.
\end{proof}

This is the entire model implemented in \texttt{stern\_model.py}. Two
sanity checks confirm it behaves correctly:
\begin{itemize}
\item As $\tau \to 0$ (end of game) with $m \ne 0$: $\dfrac{m}{\sigma\sqrt{\tau}} \to \pm\infty$,
so the probability collapses to 0 or 1 -- correct, since the outcome is
essentially decided once the clock runs out.
\item At $\tau = 1$ (kickoff) with $m = 0$: the probability is exactly
$\Phi(0) = 0.5$ -- correct, an untied prior with no score yet should be a
coin flip \emph{if we ignore the pre-game Elo information}. In practice we
blend in the Elo prior most strongly when $\tau$ is close to 1, precisely
because the pure score-based signal is uninformative there (see
\texttt{stern\_model.py}, \texttt{blend\_weight} parameter).
\end{itemize}

\subsection{Connection to first-passage / hitting-time problems}

An equivalent way to view this result: the home team loses if and only if
$M(t)$ hits the barrier at $0$ from above before the clock expires (once
trailing, ties are possible but let's set aside the discreteness). This
frames the win probability as a \emph{hitting probability} for Brownian
motion against a fixed barrier, a standard object in stochastic calculus
that also underlies barrier option pricing in the Black-Scholes framework.
The techniques are the same family used to price knock-in/knock-out
options; this project is a discrete-outcome, sports-flavored version of
the exact same mathematics.

\section{Calibrating $\sigma$}

$\sigma$ is the only free parameter in the model, and it must be estimated
from historical data before the model can be used. Rearranging the model
assumption:
\[
Y := M_{\text{final}} - M(t_0) \;\sim\; \mathcal{N}(0,\ \sigma^2 \tau)
\quad\Longrightarrow\quad
\frac{Y}{\sqrt{\tau}} \sim \mathcal{N}(0, \sigma^2)
\]
regardless of $\tau$. So, pooling every (game, snapshot) pair across
historical data, the sample variance of $Y/\sqrt{\tau}$ is a consistent
estimator of $\sigma^2$:
\[
\hat{\sigma}^2 = \widehat{\mathrm{Var}}\!\left(\frac{Y}{\sqrt{\tau}}\right)
= \frac{1}{n-1}\sum_{i=1}^n \left(\frac{y_i}{\sqrt{\tau_i}} - \bar{y/\sqrt{\tau}}\right)^2.
\]
This is simple method-of-moments estimation (implemented in
\texttt{fit\_sigma.py}) -- no maximum likelihood machinery is required
because the model assumption directly gives us a quantity with constant
variance to estimate.

\section{Measuring model quality: calibration and the Brier score}

A probabilistic forecast is \textbf{calibrated} if, among all the times the
model says ``probability $p$'', the event actually happens a fraction $p$
of the time. Formally, binning predictions into buckets $B_k$ and checking
\[
\frac{1}{|B_k|}\sum_{i \in B_k} \mathbf{1}[\text{home won}_i] \;\approx\; \bar{p}_{B_k}
\]
for every bucket $k$ is exactly the reliability diagram produced in
\texttt{calibration.py}.

A complementary scalar summary is the \textbf{Brier score}, the mean
squared error between predicted probability and the binary outcome:
\[
\text{Brier} = \frac{1}{n}\sum_{i=1}^n (p_i - o_i)^2, \qquad o_i \in \{0,1\}.
\]
It is minimized (in expectation) exactly when $p_i$ equals the true
probability of the event, which is why it is a \emph{proper scoring rule}:
a forecaster cannot improve their expected score by reporting anything
other than their honest belief. A constant 0.5 forecast scores 0.25; our
backtest on simulated data (Section 7) scores well below that.

\section{From edge to position size: the Kelly criterion}

Suppose the model believes $\mathbb{P}(E) = p$, but the market prices a YES
contract at $q$ dollars (so a NO-side profit of $q$ per contract if $E$
does not occur, and a YES-side profit of $1-q$ per contract if it does).
If we stake a fraction $f$ of our bankroll $W$ on YES, our bankroll after
the outcome is either
\[
W(1 + f\cdot b) \quad \text{(win, probability } p), \qquad
W(1 - f) \quad \text{(lose, probability } 1-p)
\]
where $b = \dfrac{1-q}{q}$ is the net odds per dollar staked. Kelly's
criterion chooses $f$ to maximize the expected log bankroll growth rate:
\[
g(f) = p\ln(1+fb) + (1-p)\ln(1-f).
\]
Differentiating and setting $g'(f) = 0$:
\[
g'(f) = \frac{pb}{1+fb} - \frac{1-p}{1-f} = 0
\quad\Longrightarrow\quad
f^\ast = \frac{pb - (1-p)}{b} = \frac{p(1+b)-1}{b}.
\]
This is the formula implemented in \texttt{kelly.py}. Two practical notes:

\begin{itemize}
\item \textbf{Fractional Kelly.} $f^\ast$ assumes $p$ is known exactly.
Since $p$ is really $p_{\text{model}}$, a noisy estimate, we stake only a
fraction (e.g. one quarter) of $f^\ast$ to reduce the damage from model
error -- a standard, well-justified adjustment in practice.
\item \textbf{Slippage.} The formula above uses the quoted price $q$, but
real fills are worse than the quote for any nonzero size. \texttt{kelly.py}
haircuts $q$ by a fixed number of cents before computing $f^\ast$ as a
simple, conservative correction.
\end{itemize}

\section{Empirical validation on simulated data}

Before trusting this pipeline on real data, we validate it on data
simulated directly from the model's own assumptions (implemented in
\texttt{generate\_mock\_data.py}): margin paths are literally simulated as
Brownian motion with a known $\sigma = 13.5$. This is a standard sanity
check in quantitative work: if the estimation and backtesting code cannot
recover known ground truth from data generated by the model itself,
there is a bug, full stop.

\textbf{Result.} Fitting $\sigma$ on $300$ simulated games ($4{,}500$
snapshots) recovered $\hat\sigma = 13.54$, against a true value of $13.5$.
The resulting calibration plot (Figure~\ref{fig:calib}) shows predicted and
realized win rates closely tracking the 45-degree line across all
probability buckets, and the backtest Brier score was $0.157$, well below
the $0.25$ coin-flip baseline.

\begin{figure}[h]
\centering
\includegraphics[width=0.65\textwidth]{../output/calibration_plot.png}
\caption{Reliability diagram on simulated data. Point size reflects the
number of observations in each probability bucket.}
\label{fig:calib}
\end{figure}

This confirms the pipeline is implemented correctly. The next step (see
README) is to re-run the same pipeline on real historical play-by-play
data, where $\hat\sigma$ and the calibration plot will reflect genuine
model error rather than pure estimation noise -- that gap between
``clean-data calibration'' and ``real-data calibration'' is itself an
informative measurement of how good the Brownian-motion approximation
actually is for real games.

\section{Limitations and honest caveats}

\begin{itemize}
\item \textbf{Constant $\sigma$.} Real scoring volatility is not constant
across a game -- end-of-game situations (intentional fouling in
basketball, clock management in football) systematically distort the
simple driftless-Brownian-motion picture. A natural extension is to fit
$\sigma(\tau)$ as a function of time remaining rather than a single
constant.
\item \textbf{No drift term.} We assumed zero drift in the in-game model,
folding all pre-game strength information into the Elo prior instead.
An alternative model could carry a drift term through the whole game.
\item \textbf{Discreteness.} Brownian motion is continuous, but football
and basketball scores jump in increments of 2, 3, 6, 7, or 8 points. The
approximation is reasonable away from the very end of the game but breaks
down near a tied score with little time left, where exact possession-based
models (Markov chains over game states) do better.
\item \textbf{Market microstructure.} We compare our model probability to
Kalshi's last traded price, but real execution faces a bid-ask spread and
thin order books; the Kelly sizing in \texttt{kelly.py} includes only a
simple slippage haircut, not a full market-impact model.
\end{itemize}

\end{document}
